\documentclass[12pt,a4paper]{article}

\usepackage[utf8]{inputenc}
\usepackage[T1]{fontenc}
\usepackage[english]{babel}
\usepackage{geometry}
\usepackage{graphicx}
\usepackage{booktabs}
\usepackage{array}
\usepackage{longtable}
\usepackage{listings}
\usepackage{xcolor}
\usepackage{hyperref}
\usepackage{fancyhdr}

\geometry{
    left=2.5cm,
    right=2.5cm,
    top=2.5cm,
    bottom=2.5cm
}

\setlength{\headheight}{15pt}

\pagestyle{fancy}
\fancyhf{}
\fancyhead[L]{Hotel Booking Manager}
\fancyhead[R]{DBMS Term Project}
\fancyfoot[C]{\thepage}

\hypersetup{
    colorlinks=true,
    linkcolor=blue,
    urlcolor=blue
}

\lstset{
    basicstyle=\ttfamily\small,
    breaklines=true,
    frame=single,
    columns=fullflexible
}

\title{
    \textbf{Hotel Booking Manager}\\
    \large Database Management Systems Term Project
}

\author{Abdelillah Taleb}
\date{\today}

\begin{document}

\maketitle

\begin{abstract}
The Hotel Booking Manager is a database-driven application for managing
hotel rooms, guests, bookings, and additional services. The project uses
PostgreSQL as the relational database, FastAPI as the REST API layer, and
a Python Tkinter desktop application as the frontend.

PostgreSQL and the FastAPI backend are deployed using Docker Compose.
The desktop frontend communicates exclusively with the REST API and is
distributed as a Debian package. Write operations are protected using an
X-API-Key.
\end{abstract}

\tableofcontents
\newpage

\section{Introduction}

The goal of this project is to implement a complete database application
consisting of a relational database, a REST API, and a desktop frontend.

The Hotel Booking Manager provides functionality for:

\begin{itemize}
    \item displaying hotel rooms,
    \item displaying existing bookings,
    \item creating new bookings,
    \item displaying booking statistics,
    \item creating additional hotel services,
    \item preventing overlapping room bookings,
    \item protecting write operations with an API key.
\end{itemize}

The system follows a layered architecture. The frontend does not connect
directly to PostgreSQL. All database access is performed through the
FastAPI REST interface.

\section{System Architecture}

The application consists of three main components:

\begin{enumerate}
    \item PostgreSQL database,
    \item FastAPI REST backend,
    \item Tkinter desktop frontend.
\end{enumerate}

PostgreSQL and FastAPI run in separate Docker containers and are
orchestrated using Docker Compose.

The overall communication flow is:

\begin{center}
\texttt{Tkinter Frontend $\rightarrow$ FastAPI REST API $\rightarrow$ PostgreSQL}
\end{center}

The frontend sends HTTP requests to the API. The API validates the
requests, executes the required SQL operations, and returns JSON
responses.

\section{Database Design}

\subsection{Tables}

The database contains six tables:

\begin{itemize}
    \item \texttt{guest}
    \item \texttt{room\_category}
    \item \texttt{room}
    \item \texttt{booking}
    \item \texttt{service}
    \item \texttt{booking\_service}
\end{itemize}

The \texttt{booking\_service} table implements an N:M relationship
between bookings and services.

\subsection{Guest}

The \texttt{guest} table stores guest information.

Important attributes include:

\begin{itemize}
    \item \texttt{guest\_id}
    \item \texttt{first\_name}
    \item \texttt{last\_name}
    \item \texttt{email}
    \item \texttt{phone}
\end{itemize}

\subsection{Room Category}

The \texttt{room\_category} table stores information shared by rooms of
the same category.

Important attributes include:

\begin{itemize}
    \item category name,
    \item capacity,
    \item price per night.
\end{itemize}

Examples of room categories are Single, Double, and Suite.

\subsection{Room}

The \texttt{room} table stores individual hotel rooms.

Each room references one room category using a foreign key.

Important attributes include:

\begin{itemize}
    \item room number,
    \item floor,
    \item active status,
    \item category ID.
\end{itemize}

\subsection{Booking}

The \texttt{booking} table stores reservations.

Important attributes include:

\begin{itemize}
    \item check-in date,
    \item check-out date,
    \item booking status,
    \item guest ID,
    \item room ID.
\end{itemize}

A database constraint ensures that the check-out date must be later than
the check-in date.

\subsection{Service}

The \texttt{service} table stores optional hotel services such as
breakfast, parking, or late check-out.

Each service has a unique name and a non-negative price.

\subsection{Booking Service}

The \texttt{booking\_service} table connects bookings and services.

Its composite primary key consists of:

\begin{itemize}
    \item \texttt{booking\_id}
    \item \texttt{service\_id}
\end{itemize}

The table also stores the service quantity.

\section{Normalization}

The relational model is normalized to Third Normal Form (3NF).

The design avoids unnecessary duplication by separating independent
entities into individual tables.

For example, room category information such as capacity and price is
stored once in \texttt{room\_category} instead of being repeated for
every room.

The N:M relationship between bookings and services is represented by the
junction table \texttt{booking\_service}.

Each non-key attribute depends on the key, the whole key, and nothing but
the key.

\section{Business Rules}

\subsection{Booking Date Validation}

A booking is only valid if:

\begin{center}
\texttt{check\_out > check\_in}
\end{center}

Invalid date ranges are rejected by the API with HTTP status code 400.

\subsection{Prevention of Overlapping Bookings}

The database uses a PostgreSQL exclusion constraint to prevent two active
bookings from overlapping for the same room.

The constraint uses a date range and applies to bookings with status
\texttt{confirmed} or \texttt{checked\_in}.

An overlapping booking request is returned by the API as:

\begin{center}
\texttt{409 Conflict}
\end{center}

This rule is implemented at database level so that data consistency does
not depend only on application code.

\subsection{Service Validation}

Service prices must be non-negative.

Service names must be unique.

A duplicate service name results in HTTP status code 409.

\section{REST API}

The REST API is implemented with FastAPI.

The main endpoints are shown in Table~\ref{tab:endpoints}.

\begin{table}[h]
\centering
\begin{tabular}{lll}
\toprule
Method & Endpoint & Description \\
\midrule
GET & \texttt{/rooms} & List all rooms \\
GET & \texttt{/bookings} & List all bookings \\
GET & \texttt{/statistics/bookings} & Booking statistics \\
POST & \texttt{/bookings} & Create a booking \\
POST & \texttt{/services} & Create a service \\
\bottomrule
\end{tabular}
\caption{Main REST API endpoints}
\label{tab:endpoints}
\end{table}

A root endpoint is also available at:

\begin{center}
\texttt{/}
\end{center}

It can be used as a simple API health check.

FastAPI automatically provides Swagger documentation at:

\begin{center}
\texttt{http://localhost:8003/docs}
\end{center}

\subsection{JOIN Queries}

The API uses SQL JOIN operations to combine related information.

For example, the room endpoint joins the \texttt{room} table with
\texttt{room\_category} so that category name, capacity, and price can be
returned together with room information.

The bookings endpoint joins:

\begin{itemize}
    \item \texttt{booking},
    \item \texttt{guest},
    \item \texttt{room}.
\end{itemize}

This allows the frontend to display guest and room information without
performing multiple API requests.

\subsection{Aggregate Query}

The endpoint

\begin{center}
\texttt{/statistics/bookings}
\end{center}

uses PostgreSQL aggregation to count bookings by status.

The response contains values for:

\begin{itemize}
    \item total bookings,
    \item confirmed bookings,
    \item checked-in bookings,
    \item checked-out bookings,
    \item cancelled bookings.
\end{itemize}

\section{API Authentication}

Write operations are protected with an \texttt{X-API-Key} HTTP header.

The protected endpoints are:

\begin{itemize}
    \item \texttt{POST /bookings}
    \item \texttt{POST /services}
\end{itemize}

The API key is read from an environment variable:

\begin{lstlisting}
API_KEY = os.getenv("API_KEY", "hotel-booking-key")
\end{lstlisting}

The real deployment value is stored in the local \texttt{.env} file and
is not committed to Git.

A valid request contains a header similar to:

\begin{lstlisting}
X-API-Key: <your-api-key>
\end{lstlisting}

Requests without a valid API key return:

\begin{center}
\texttt{401 Unauthorized}
\end{center}

\section{Environment Configuration}

Deployment credentials and secrets are stored in a \texttt{.env} file.

An example configuration is:

\begin{lstlisting}
POSTGRES_DB=hotel_booking_ataleb
POSTGRES_USER=hotel_booking
POSTGRES_PASSWORD=<database-password>
API_KEY=<api-key>
\end{lstlisting}

The \texttt{.env} file is included in \texttt{.gitignore} and therefore
is not stored in the public repository.

The FastAPI backend reads its database configuration using environment
variables such as:

\begin{itemize}
    \item \texttt{DB\_HOST}
    \item \texttt{DB\_PORT}
    \item \texttt{DB\_NAME}
    \item \texttt{DB\_USER}
    \item \texttt{DB\_PASSWORD}
\end{itemize}

\section{Docker Deployment}

PostgreSQL and FastAPI are deployed using Docker Compose.

The project contains:

\begin{itemize}
    \item \texttt{compose.yaml}
    \item \texttt{backend/Dockerfile}
    \item \texttt{backend/requirements.txt}
\end{itemize}

The system can be started from the project directory with:

\begin{lstlisting}
docker compose up -d --build
\end{lstlisting}

The status of the containers can be checked with:

\begin{lstlisting}
docker compose ps
\end{lstlisting}

The PostgreSQL container uses a health check based on
\texttt{pg\_isready}.

The API container starts after PostgreSQL becomes healthy.

The database schema and seed data are initialized using:

\begin{itemize}
    \item \texttt{database/schema.sql}
    \item \texttt{database/seed.sql}
\end{itemize}

The API is exposed on port 8003.

The containers can be stopped using:

\begin{lstlisting}
docker compose down
\end{lstlisting}

\section{Desktop Frontend}

The frontend is implemented in Python using Tkinter.

At startup, a connection dialog asks the user for:

\begin{itemize}
    \item API URL,
    \item X-API-Key.
\end{itemize}

The default development API URL is:

\begin{center}
\texttt{http://localhost:8003}
\end{center}

The entered API key is used for protected POST requests.

The frontend provides the following views:

\begin{itemize}
    \item room overview,
    \item booking overview,
    \item booking creation form,
    \item booking statistics,
    \item service creation form.
\end{itemize}

All communication is performed using HTTP requests to FastAPI.
The frontend never accesses PostgreSQL directly.

\section{Debian Package}

The frontend is distributed as a Debian package.

The generated package is:

\begin{center}
\texttt{dist/hotel-booking-manager\_1.0.0\_all.deb}
\end{center}

The package contains:

\begin{itemize}
    \item the Python frontend application,
    \item a launcher command,
    \item Debian package metadata.
\end{itemize}

The launcher is installed as:

\begin{center}
\texttt{/usr/bin/hotel-booking-manager}
\end{center}

The application source is installed as:

\begin{center}
\texttt{/usr/share/hotel-booking-manager/main.py}
\end{center}

The package depends on:

\begin{itemize}
    \item \texttt{python3},
    \item \texttt{python3-tk}.
\end{itemize}

The package can be built with:

\begin{lstlisting}
dpkg-deb --build \
  packaging/hotel-booking-manager \
  dist/hotel-booking-manager_1.0.0_all.deb
\end{lstlisting}

It can be inspected using:

\begin{lstlisting}
dpkg-deb --info \
  dist/hotel-booking-manager_1.0.0_all.deb

dpkg-deb --contents \
  dist/hotel-booking-manager_1.0.0_all.deb
\end{lstlisting}

On a Debian machine with administrator privileges, the package can be
installed with:

\begin{lstlisting}
sudo dpkg -i \
  dist/hotel-booking-manager_1.0.0_all.deb
\end{lstlisting}

After installation, the frontend can be started with:

\begin{lstlisting}
hotel-booking-manager
\end{lstlisting}

The lecture server account used during development does not provide
sudo privileges. Therefore, package creation and inspection can be
performed by the project user, while system-wide installation requires
an account with administrative permissions.

\section{Testing}

The implementation was tested at database, API, deployment, and package
levels.

\subsection{Database Tests}

The database constraint for overlapping bookings was tested directly.

An attempt to create an overlapping booking for the same room was
rejected successfully.

\subsection{API Tests}

The following cases were tested:

\begin{itemize}
    \item GET request for rooms,
    \item GET request for bookings,
    \item booking statistics,
    \item valid booking creation,
    \item invalid booking date range,
    \item overlapping booking,
    \item invalid guest or room reference,
    \item valid service creation,
    \item duplicate service name,
    \item missing API key,
    \item valid API key.
\end{itemize}

A missing API key returned:

\begin{center}
\texttt{401 Unauthorized}
\end{center}

A valid API key stored in the environment configuration was accepted and
a service was created with HTTP status code 201.

\subsection{Docker Tests}

Docker Compose was tested successfully.

The PostgreSQL container reached the healthy state and the FastAPI
container started correctly.

A request to:

\begin{center}
\texttt{GET /rooms}
\end{center}

returned the seeded room data from the PostgreSQL container.

\subsection{Debian Package Tests}

The Debian package was successfully built using \texttt{dpkg-deb}.

Its metadata and file contents were inspected.

The package contains both:

\begin{itemize}
    \item the executable launcher,
    \item the Tkinter frontend source.
\end{itemize}

The package was also extracted to a temporary directory without errors.

\section{Project Structure}

The relevant project structure is:

\begin{lstlisting}
DBMS_10/
|-- backend/
|   |-- Dockerfile
|   |-- main.py
|   `-- requirements.txt
|-- database/
|   |-- schema.sql
|   `-- seed.sql
|-- frontend/
|   `-- main.py
|-- packaging/
|   `-- hotel-booking-manager/
|-- dist/
|   `-- hotel-booking-manager_1.0.0_all.deb
|-- example-documentation/
|   `-- documentation.tex
|-- compose.yaml
|-- Makefile
|-- README.md
`-- .gitignore
\end{lstlisting}

\section{Version Control}

The project source code is maintained using Git.

The repository contains the application code, SQL schema, seed data,
Docker configuration, Debian package files, README, and LaTeX
documentation.

Secrets such as the database password and API key are excluded from
version control using the \texttt{.gitignore} file.

\section{Conclusion}

The Hotel Booking Manager demonstrates the complete implementation of a
small database application.

The project combines:

\begin{itemize}
    \item relational database design,
    \item Third Normal Form,
    \item primary and foreign keys,
    \item N:M relationships,
    \item SQL constraints,
    \item JOIN queries,
    \item aggregate queries,
    \item REST API development,
    \item API authentication,
    \item Docker Compose deployment,
    \item desktop frontend development,
    \item Debian packaging.
\end{itemize}

The database is protected from inconsistent booking data using database
constraints, while FastAPI provides a controlled interface between the
frontend and PostgreSQL.

Docker Compose provides reproducible deployment of the database and API,
and the Debian package provides a standard installation format for the
desktop frontend.

\end{document}
